\section{The space of eigenvalue responses}\label{sec:responses}

Fix a size $n\in\mathcal S$, a generic input pair $P$, and a splitting
field $\Omega_P$ of $W(P)$ over $\operatorname{Frac}K[[P]]$.
Write its distinct roots as $\lambda_h$, with spectral projections
$\pi_h$, $1\le h\le e_n$.  Each projection has rank $t$.
We also use these data after embedding into the valued fields of
Section~\ref{sec:formal-generic}.

We seek the possible leading changes of the eigenvalues at $P$.
Besides varying $P$ itself, we allow an excursion through another
input block and a return to $P$.  Their linear span will determine
which eigenvalues can be varied independently.

\subsection{Loop and paired responses}\label{subsec:response-formulas}

Upper triangular block evaluations are also the basic mechanism of
noncommutative difference-differential calculus; see
\cite[Proposition~2.2 and Theorem~3.11]{KaliuzhnyiVerbovetskyiVinnikov2014}.
Here the admissible evaluations of Section~\ref{sec:formal-generic}
let us extract scalar eigenvalue responses from those block calculations.

For a perturbation $\mathcal W$ of $W(P)$, put $\pi=\pi_h$,
$\bar\pi=1-\pi$, and
\[
 C_h(\mathcal W)=\bar\pi(\mathcal W-\lambda_h)\bar\pi.
\]
This operator acts on the complementary space. Its constant term is
invertible, so its inverse has a unique expansion in either commuting
or free formal parameters. Define the fixed-root Schur expression by
\begin{equation}
S_h(\mathcal W)=
 \pi(\mathcal W-\lambda_h)\pi
 -\pi\mathcal W\bar\pi
  C_h(\mathcal W)^{-1}
  \bar\pi\mathcal W\pi.
 \label{eq:fixed-schur}
\end{equation}
This expression records the obstruction to keeping the root $\lambda_h$
fixed. For $u$ in the image of $\pi$, the vector
\[
 u-C_h(\mathcal W)^{-1}\bar\pi\mathcal W\pi u
\]
is the unique vector projecting to $u$ whose complementary component
under $\mathcal W-\lambda_h$ vanishes. Its remaining component is
$S_h(\mathcal W)u$. Thus $S_h(\mathcal W)=0$ means that the fixed-root
kernel is a graph over the original eigenspace.

When the constant output is split semisimple and these equations are
imposed for every one of its distinct roots, the kernel graphs combine
into a formal change of basis. The output remains conjugate to its
constant term and satisfies that constant term's minimal polynomial.
Lemma~\ref{lem:schur-polynomial} proves the equivalence with the
polynomial identity at every finite word order.

The same Schur expression is defined near $\operatorname{diag}(P,C_0)$
if $\lambda_h$ is absent from $W(C_0)$, embedding $\pi_h$ at the $P$
block. In this use, equations indexed by labels at $P$ record those
fixed roots only.

For a loop direction $Y=(Y_1,Y_2)$ at $P$, define the scalar
$D_h(P;Y)$ by
\begin{equation}
\pi_h W'_P(Y)\pi_h=D_h(P;Y)\pi_h.
 \label{eq:loop-response}
\end{equation}
This compression is scalar: apply Lemma~\ref{lem:scalar-branches}
to the commuting deformation $P+xY$ and take its first coefficient.
Equivalently, $D_h(P;Y)$ is the differential of the branch
$\lambda_h$ in direction $Y$.  By
Lemma~\ref{lem:nonconstant-roots}, each of these linear functionals is
nonzero.  Because $K$ is infinite, we can choose a constant direction
$v$ over $K$ with
\begin{equation}
D_h(P;v)\ne0\qquad(1\le h\le e_n).
 \label{eq:nonzero-direction}
\end{equation}
Indeed the forbidden directions form a finite union of proper
$K$-linear subspaces, even though their defining coefficients belong
to $\Omega_P$.

Let $Q$ be an independent generic point of size $n'\in\mathcal S$.
Let $U$ be an input arrow $P\leftarrow Q$ and $V$ an input arrow
$Q\leftarrow P$; each is a pair of rectangular scalar matrices.
Form the triangular input pair
\[
 T=\begin{pmatrix}P&U&0\\0&Q&V\\0&0&P\end{pmatrix}.
\]
Set
\begin{equation}
M_h(P,Q;U,V)=\pi_h\left(
 W(T)_{13}
 -W(T)_{12}(W(Q)-\lambda_h)^{-1}W(T)_{23}
 \right)\pi_h.
 \label{eq:paired-response}
\end{equation}
Block paths show that this is bilinear in $U,V$.  The inverse exists
because the spectra at independent points are disjoint.

\begin{lemma}\label{lem:paired-scalar}
There is a scalar $m_h(P,Q;U,V)$ such that
\[
 M_h(P,Q;U,V)=m_h(P,Q;U,V)\pi_h.
\]
It is the coefficient of $xy$ in the branch reducing to $\lambda_h$
for the commuting input deformation
\[
 \begin{pmatrix}P&xU\\yV&Q\end{pmatrix}.
\]
The same assertion holds at an admissible specialized pair whose
spectra are disjoint and whose combined output is split semisimple
with the maximal number $e_{n+n'}$ of distinct roots.
\end{lemma}

\begin{proof}
At $x=y=0$, disjointness and Lemma~\ref{lem:stable-multiplicity}
give $e_n+e_{n'}=e_{n+n'}$ roots.  Thus
Lemma~\ref{lem:scalar-branches} applies.  Either arrow alone leaves
the spectrum unchanged, so the chosen branch has no terms depending
on $x$ alone or on $y$ alone.

Compute the $xy$ coefficient of its fixed-root Schur expression.
The direct return from $P$ through $Q$ gives $W(T)_{13}$.  Eliminating
the $Q$ block gives the second term in
\eqref{eq:paired-response}.  The other eigenspaces in the $P$ block
give no further term of this order: their couplings to $\pi_h$ have
zero constant term and no first-order term under these off-diagonal
input perturbations.  Replacing the fixed root by the scalar branch
subtracts its $xy$ coefficient times $\pi_h$.  The resulting Schur
expression vanishes, because the branch eigenspace projects
isomorphically to the base eigenspace.  This proves the formula.
The specialized assertion uses exactly the same argument and the
specialization-stable degree bound.
\end{proof}

\subsection{Coefficient spans and their specializations}
\label{subsec:active-space}

At $Q=0$ the matrices $W(Q)-\lambda_h$ are invertible, since no
$\lambda_h$ belongs to the spectrum of $w(0)$.  Hence
\eqref{eq:paired-response} has a formal expansion in the entries of
$Q$ over $\Omega_P$.  Define
\[
 m_h=\frac1t\operatorname{tr}M_h,
 \qquad m=(m_h)_{h=1}^{e_n},\qquad D=(D_h)_{h=1}^{e_n}.
\]
Here $D=D(P;\cdot):M_n(\Omega_P)^2\to\Omega_P^{e_n}$ is the linear
map of loop responses. Its columns are its values on the standard
coordinate directions in the two input matrices.
The scalar identity of Lemma~\ref{lem:paired-scalar} holds also for
these formal expansions, by clearing their determinants and applying
Lemma~\ref{lem:generic-detection}.

Define $\mathcal T_P\subseteq\Omega_P^{e_n}$ to be the linear span
of the columns of $D$ and of every coefficient vector of $m$, for
every allowed partner size $n'\in\mathcal S$ and every pair of
coordinate arrows $U,V$.  Write
\[
 a_n=\dim_{\Omega_P}\mathcal T_P.
\]
Although the definition permits infinitely many partner sizes and
coefficients, the ambient space is finite dimensional.

\begin{lemma}[Specialization and finite witnesses]
\label{lem:active-specialization}
The scalar identities for $M_h$ and every linear identity expressing
$m\in\mathcal T_P$ hold after any admissible specialization of the
partner input for which the required inverses exist.  For identities
on a subset of labels, only the inverses for those labels are needed.

After a field extension, the columns of $D$ and finitely many paired
responses span $\mathcal T_P$.  These responses can be taken at
independent small generic partner points, whose sizes need not be
equal.
\end{lemma}

\begin{proof}
For the first assertion, apply any $\Omega_P$-linear functional
annihilating $\mathcal T_P$ to the formal tuple $m$.  The result is
zero.  Multiply by the finitely many determinants of
$W(Q)-\lambda_h$ that occur.  The cleared expression is a finite
sum of products of base-point scalars and matrix-entry series; its
auxiliary coefficients have a uniform valuation bound.  Its formal
identity is therefore equivalent to its identity at a small generic
partner, by Lemma~\ref{lem:generic-detection}.  It persists under the
admissible specializations of Lemma~\ref{lem:admissible-evaluation}.
Divide by the specialized determinants when they are nonzero.  The
same argument, entry by entry, proves the scalar identities for
$M_h$.  It also proves the projected assertion, without introducing
any spectral labels for the partner input.

Choose finitely many coefficient vectors, together with columns of
$D$, that span $\mathcal T_P$, and select a nonzero minor witnessing
their rank.  For each column requiring a partner, use a separate set
of partner variables.  By multilinearity, the corresponding
determinant of the series columns has that nonzero minor as the
coefficient of a specified product of monomials.  It is therefore a
nonzero formal series.  Clear the column denominators.  The numerator
is still nonzero, and its coefficients are uniformly bounded below
in valuation.  Independent small generic evaluations keep this
numerator and its denominators nonzero, by
Lemma~\ref{lem:generic-detection}.  The resulting values attain the
full dimension.  The first assertion places all of them in the
extended active space.

This argument evaluates cleared numerators, not the possibly
divergent geometric expansions of their inverses.  The distinction
is precisely the one made in Section~\ref{subsec:formal-conventions}.
\end{proof}

The lemma also applies after adjoining Taylor or Laurent parameters.
One first clears determinants and then works coefficientwise as in
Lemma~\ref{lem:admissible-evaluation}.  In particular, it permits
triangular partners with repeated base blocks, provided the inverses
required for the selected target roots exist in the chosen Laurent
field.

\subsection{A partition of the eigenvalue labels}
\label{subsec:active-partition}

The active space has a stronger structure than a general linear
subspace.  Its coordinates form groups whose responses are always in
a fixed ratio.

\begin{proposition}\label{prop:active-partition}
For the direction $v$ chosen in \eqref{eq:nonzero-direction}, there
is a partition of $\{1,\ldots,e_n\}$ such that
\[
 \mathcal T_P=
 \left\{c\in\Omega_P^{e_n}:
 \frac{c_h}{D_h(P;v)}\text{ is constant on each part}\right\}.
\]
\end{proposition}

\begin{proof}
We first prove that $\mathcal T_P$ is closed under the operation
\begin{equation}
(a,b)\longmapsto
 \left(\frac{a_hb_h}{D_h(P;v)}\right)_h.
 \label{eq:weighted-product}
\end{equation}
Choose two finite spanning sets supplied by
Lemma~\ref{lem:active-specialization}, using independent partner
points for the two sets.  It suffices to prove the assertion for a
column $a$ in the first set and a column $b$ in the second, over their
joint coefficient field $E$.

Put $P'=P+\xi v$.  Use a triangular chain with endpoints $P,P$,
middle point $P'$, an optional point $Q_1$ between the first endpoint
and $P'$, and an optional point $Q_2$ between $P'$ and the last
endpoint.  A direct step represents a loop column; an excursion
through $Q_i$ represents a paired column.  Put the arrows defining
$a$ and $b$ on the successive steps, with no other input arrows.
Group the internal positions as a single partner $C(\xi)$.  Its
size belongs to $\mathcal S$.

For every target label, $W(C(\xi))-\lambda_h$ is invertible over
$E((\xi))$.  The optional blocks have spectra disjoint from $P$.
At the middle block, the branch reducing to $\lambda_h$ moves with
nonzero derivative $D_h(P;v)$; all other branches are already away
from $\lambda_h$.  Thus
\begin{equation}
\left.\xi\bigl(W(P')-\lambda_h\bigr)^{-1}\right|_{\xi=0}
   =\frac{\pi_h}{D_h(P;v)}.
 \label{eq:middle-residue}
\end{equation}
This follows directly by scalar branch separation.

By Lemma~\ref{lem:active-specialization}, the paired response through
$C(\xi)$ belongs to $\mathcal T_P\otimes E((\xi))$.  Eliminate
the optional blocks first in its Schur formula.  Triangularity means
that their elimination does not change the middle diagonal block.
All terms are regular at $\xi=0$ except the middle inverse in
\eqref{eq:middle-residue}.  The effective first-to-middle coupling,
compressed by $\pi_h$ at both ends and evaluated at zero, is
$a_h\pi_h$.  For a direct step this is
\eqref{eq:loop-response}; for an excursion it is
\eqref{eq:paired-response}.  The corresponding middle-to-last
coupling is $b_h\pi_h$.

The singular term occurs with the minus sign in the Schur
complement.  Multiplying by $\xi$ and reducing at zero therefore
gives
\[
 \left(-\frac{a_hb_h}{D_h(P;v)}\right)_h
       \in\mathcal T_P\otimes E.
\]
Taking a Laurent coefficient preserves membership in this fixed
finite dimensional space: apply its annihilating linear functionals.
Bilinearity and the two spanning sets prove closure under
\eqref{eq:weighted-product} after extension to $E$, and faithful
scalar extension proves it over $\Omega_P$.

Divide coordinates by the nowhere-zero vector $D(P;v)$.  The resulting
space $\mathcal U\subseteq\Omega_P^{e_n}$ contains $1$, since
$D(P;v)\in\mathcal T_P$, and is closed under coordinatewise
multiplication.  Declare two coordinates equivalent if every element
of $\mathcal U$ has the same value on them.  By avoiding finitely
many proper linear subspaces, choose an element of $\mathcal U$
whose values distinguish all the equivalence classes.  Lagrange
interpolation in this element gives the indicator function of each
class.  These indicators span all class-constant functions.  Thus
$\mathcal U$ is exactly that algebra of functions, proving the
proposition.
\end{proof}

We call two labels in the same part \emph{linked}.  The number of
parts is $a_n$.  Selecting one representative from each part
identifies $\mathcal T_P$ with the full coordinate space on those
representatives.  The next stage must show that these response
relations control higher-order free-letter errors, rather than only
the two elementary perturbations used to define them.
